% --- Archon physics formalization source begin ---
% archon:physics
% archon:covers PhyXMiniProblems/problem_phyx_mini_0950.lean
% archon:source-report reports/phyx_mini/problem_phyx_mini_0950.source.json

\chapter{Physics problem phyx\_mini\_0950}
\label{ch:PhyXMiniProblems_problem_phyx_mini_0950}

\paragraph{Problem source.}
\#\# Physical scenario

A voice coil in a loudspeaker has 50 turns of wire and a diameter of \textbackslash{}( 1.56\textbackslash{},\textbackslash{}text\{cm\} \textbackslash{}), and the current in the coil is \textbackslash{}( 0.950\textbackslash{},\textbackslash{}text\{A\} \textbackslash{}). Assume that the magnetic field at each point of the coil has a constant magnitude of \textbackslash{}( 0.220\textbackslash{},\textbackslash{}text\{T\} \textbackslash{}) and is directed at an angle of \textbackslash{}( 60.0\textasciicircum{}\textbackslash{}circ \textbackslash{}) outward from the normal to the plane of the coil. Let the axis of the coil be in the \textbackslash{}( y \textbackslash{})-direction. The current in the coil is in the direction shown (counterclockwise as viewed from a point above the coil on the \textbackslash{}( y \textbackslash{})-axis).

\#\# Question

Calculate the magnitude and direction of the net magnetic force on the coil.

\#\# Answer choices

- A: A: 0.454N

- B: B: -0.444N

- C: C: 0.434N

- D: D: -0.456N

\#\# Figure caption (auxiliary; use the image as primary evidence)

The image depicts a diagram of a solenoid represented by a coil wrapped around a cylindrical object on an x-y coordinate plane. 

Key elements include:

- **Cylinder:** A vertical cylindrical object is positioned along the y-axis.
  
- **Coil:** Multiple loops of coil wrap around the cylinder, depicted as several concentric circles from the side view. An arrow on the coil indicates the direction of current, labeled as "I".

- **Magnetic Fields (\textbackslash{}( \textbackslash{}vec\{B\} \textbackslash{})):** Two blue arrows with labels indicating magnetic field directions emerge symmetrically from the sides of the coil at a 60.0° angle relative to the horizontal. These arrows are labeled with vector symbols \textbackslash{}(\textbackslash{}vec\{B\}\textbackslash{}).

- **Angles:** The two angles between the magnetic field lines and the vertical axis are marked as 60.0°.

- **Axes:** The coordinate axes are marked as x and y.

Overall, the illustration demonstrates the interaction of electric current in a coil with the generated magnetic fields.

\#\# Dataset metadata

Magnetism; Physical Model Grounding Reasoning, Multi-Formula Reasoning

\paragraph{Recorded answer/context.}
B: B: -0.444N

\paragraph{Figure/image path.}
/root/proposal\_for\_physic/science-mango/phyx\_mini\_run/phyx\_data/test\_image/950.png

\paragraph{Formalization target.}
create a compiling Lean file with sorry bodies at `PhyXMiniProblems/problem\_phyx\_mini\_0950.lean`. The Lean declarations must preserve the physical quantities, dimensions or dimensional roles, figure labels, governing-law hypotheses, and final relation expressed by this problem.
Use Mathlib/Physlib names found through LeanExplore where available. If a physics API is missing, introduce faithful local abstractions rather than scalar placeholder aliases.

\begin{theorem}[Physics formalization target]
\label{thm:physics:phyx_mini_0950:target}
\lean{PhyXMiniProblems.ProblemPhyXMini0950.problem_phyx_mini_0950}
\uses{def:physics:phyx-mini-0950:phyxminiproblems-problemphyxmini0950-forcevectorinnewtons, def:physics:phyx-mini-0950:phyxminiproblems-problemphyxmini0950-forcemagnitudeinnewtons, def:physics:phyx-mini-0950:phyxminiproblems-problemphyxmini0950-yhat, def:physics:phyx-mini-0950:phyxminiproblems-problemphyxmini0950-voicecoilsetup, def:physics:phyx-mini-0950:phyxminiproblems-problemphyxmini0950-matchesvoicecoilscenario, def:physics:phyx-mini-0950:phyxminiproblems-problemphyxmini0950-matchesprimaryvoicecoilfigure, def:physics:phyx-mini-0950:phyxminiproblems-problemphyxmini0950-matchesproblemreadouts, def:physics:phyx-mini-0950:phyxminiproblems-problemphyxmini0950-hasphysicalvoicecoilparameters, def:physics:phyx-mini-0950:phyxminiproblems-problemphyxmini0950-satisfiescircularvoicecoilgeometry, def:physics:phyx-mini-0950:phyxminiproblems-problemphyxmini0950-modelsaxisymmetricvoicecoilmagneticfield, def:physics:phyx-mini-0950:phyxminiproblems-problemphyxmini0950-satisfiesvoicecoilmagneticlorentzforcelaw, def:physics:phyx-mini-0950:phyxminiproblems-problemphyxmini0950-exactforcemagnitudeinnewtons, def:physics:phyx-mini-0950:phyxminiproblems-problemphyxmini0950-pointsalongnegativeyaxis, def:physics:phyx-mini-0950:phyxminiproblems-problemphyxmini0950-matchesdisplayedsignedforce, def:physics:phyx-mini-0950:phyxminiproblems-problemphyxmini0950-isuniqueclosestsignedforcechoice}
The assigned autoformalize agent should translate this physics problem into Lean declarations in the covered file, with theorem and lemma proof bodies written as `by sorry`.
\end{theorem}
\begin{proof}
This is an autoformalization task, not a proof task. Produce statements that can later be proved by the physics prover without weakening the physical model.
\end{proof}

% --- Archon declaration topology begin ---
\section{Declaration topology}
\begin{definition}[Spatial Vector]
  \label{def:physics:phyx-mini-0950:phyxminiproblems-problemphyxmini0950-spatialvector}
  \lean{PhyXMiniProblems.ProblemPhyXMini0950.SpatialVector}
  Three-dimensional Euclidean vectors used for coherent-unit readouts
\end{definition}
\begin{proof}
  This is the declared model definition of Spatial Vector.
\end{proof}
\begin{definition}[electric Current Dimension]
  \label{def:physics:phyx-mini-0950:phyxminiproblems-problemphyxmini0950-electriccurrentdimension}
  \lean{PhyXMiniProblems.ProblemPhyXMini0950.electricCurrentDimension}
  Electric current has physical dimension C T⁻¹
\end{definition}
\begin{proof}
  This is the declared model definition of electric Current Dimension.
\end{proof}
\begin{definition}[magnetic Flux Density Dimension]
  \label{def:physics:phyx-mini-0950:phyxminiproblems-problemphyxmini0950-magneticfluxdensitydimension}
  \lean{PhyXMiniProblems.ProblemPhyXMini0950.magneticFluxDensityDimension}
  Magnetic flux density (tesla) has physical dimension M T⁻¹ C⁻¹
\end{definition}
\begin{proof}
  This is the declared model definition of magnetic Flux Density Dimension.
\end{proof}
\begin{definition}[force Dimension]
  \label{def:physics:phyx-mini-0950:phyxminiproblems-problemphyxmini0950-forcedimension}
  \lean{PhyXMiniProblems.ProblemPhyXMini0950.forceDimension}
  Force (newton) has physical dimension M L T⁻²
\end{definition}
\begin{proof}
  This is the declared model definition of force Dimension.
\end{proof}
\begin{definition}[Length Quantity]
  \label{def:physics:phyx-mini-0950:phyxminiproblems-problemphyxmini0950-lengthquantity}
  \lean{PhyXMiniProblems.ProblemPhyXMini0950.LengthQuantity}
  A nonnegative, unit-independent physical length
\end{definition}
\begin{proof}
  This is the declared model definition of Length Quantity.
\end{proof}
\begin{definition}[Current Magnitude Quantity]
  \label{def:physics:phyx-mini-0950:phyxminiproblems-problemphyxmini0950-currentmagnitudequantity}
  \lean{PhyXMiniProblems.ProblemPhyXMini0950.CurrentMagnitudeQuantity}
  \uses{def:physics:phyx-mini-0950:phyxminiproblems-problemphyxmini0950-electriccurrentdimension}
  A nonnegative, unit-independent electric-current magnitude
\end{definition}
\begin{proof}
  This is the declared model definition of Current Magnitude Quantity.
\end{proof}
\begin{definition}[Magnetic Flux Density Magnitude Quantity]
  \label{def:physics:phyx-mini-0950:phyxminiproblems-problemphyxmini0950-magneticfluxdensitymagnitudequantity}
  \lean{PhyXMiniProblems.ProblemPhyXMini0950.MagneticFluxDensityMagnitudeQuantity}
  \uses{def:physics:phyx-mini-0950:phyxminiproblems-problemphyxmini0950-magneticfluxdensitydimension}
  A nonnegative, unit-independent magnetic-flux-density magnitude
\end{definition}
\begin{proof}
  This is the declared model definition of Magnetic Flux Density Magnitude Quantity.
\end{proof}
\begin{definition}[Force Vector Quantity]
  \label{def:physics:phyx-mini-0950:phyxminiproblems-problemphyxmini0950-forcevectorquantity}
  \lean{PhyXMiniProblems.ProblemPhyXMini0950.ForceVectorQuantity}
  \uses{def:physics:phyx-mini-0950:phyxminiproblems-problemphyxmini0950-spatialvector, def:physics:phyx-mini-0950:phyxminiproblems-problemphyxmini0950-forcedimension}
  A signed, unit-independent spatial force vector
\end{definition}
\begin{proof}
  This is the declared model definition of Force Vector Quantity.
\end{proof}
\begin{definition}[nonnegative SIReadout]
  \label{def:physics:phyx-mini-0950:phyxminiproblems-problemphyxmini0950-nonnegativesireadout}
  \lean{PhyXMiniProblems.ProblemPhyXMini0950.nonnegativeSIReadout}
  Read a nonnegative physical quantity in coherent SI units
\end{definition}
\begin{proof}
  This is the declared model definition of nonnegative SIReadout.
\end{proof}
\begin{definition}[vector SIReadout]
  \label{def:physics:phyx-mini-0950:phyxminiproblems-problemphyxmini0950-vectorsireadout}
  \lean{PhyXMiniProblems.ProblemPhyXMini0950.vectorSIReadout}
  \uses{def:physics:phyx-mini-0950:phyxminiproblems-problemphyxmini0950-spatialvector}
  Read a dimensionful spatial vector in coherent SI components
\end{definition}
\begin{proof}
  This is the declared model definition of vector SIReadout.
\end{proof}
\begin{definition}[length In Meters]
  \label{def:physics:phyx-mini-0950:phyxminiproblems-problemphyxmini0950-lengthinmeters}
  \lean{PhyXMiniProblems.ProblemPhyXMini0950.lengthInMeters}
  \uses{def:physics:phyx-mini-0950:phyxminiproblems-problemphyxmini0950-lengthquantity, def:physics:phyx-mini-0950:phyxminiproblems-problemphyxmini0950-nonnegativesireadout}
  Metre readout of a physical length
\end{definition}
\begin{proof}
  This is the declared model definition of length In Meters.
\end{proof}
\begin{definition}[length In Centimeters]
  \label{def:physics:phyx-mini-0950:phyxminiproblems-problemphyxmini0950-lengthincentimeters}
  \lean{PhyXMiniProblems.ProblemPhyXMini0950.lengthInCentimeters}
  \uses{def:physics:phyx-mini-0950:phyxminiproblems-problemphyxmini0950-lengthquantity, def:physics:phyx-mini-0950:phyxminiproblems-problemphyxmini0950-lengthinmeters}
  Centimetre readout used by the diameter stated in the problem
\end{definition}
\begin{proof}
  This is the declared model definition of length In Centimeters.
\end{proof}
\begin{definition}[current In Amperes]
  \label{def:physics:phyx-mini-0950:phyxminiproblems-problemphyxmini0950-currentinamperes}
  \lean{PhyXMiniProblems.ProblemPhyXMini0950.currentInAmperes}
  \uses{def:physics:phyx-mini-0950:phyxminiproblems-problemphyxmini0950-currentmagnitudequantity, def:physics:phyx-mini-0950:phyxminiproblems-problemphyxmini0950-nonnegativesireadout}
  Ampere readout of the current magnitude
\end{definition}
\begin{proof}
  This is the declared model definition of current In Amperes.
\end{proof}
\begin{definition}[magnetic Flux Density In Teslas]
  \label{def:physics:phyx-mini-0950:phyxminiproblems-problemphyxmini0950-magneticfluxdensityinteslas}
  \lean{PhyXMiniProblems.ProblemPhyXMini0950.magneticFluxDensityInTeslas}
  \uses{def:physics:phyx-mini-0950:phyxminiproblems-problemphyxmini0950-magneticfluxdensitymagnitudequantity, def:physics:phyx-mini-0950:phyxminiproblems-problemphyxmini0950-nonnegativesireadout}
  Tesla readout of the magnetic-flux-density magnitude
\end{definition}
\begin{proof}
  This is the declared model definition of magnetic Flux Density In Teslas.
\end{proof}
\begin{definition}[force Vector In Newtons]
  \label{def:physics:phyx-mini-0950:phyxminiproblems-problemphyxmini0950-forcevectorinnewtons}
  \lean{PhyXMiniProblems.ProblemPhyXMini0950.forceVectorInNewtons}
  \uses{def:physics:phyx-mini-0950:phyxminiproblems-problemphyxmini0950-spatialvector, def:physics:phyx-mini-0950:phyxminiproblems-problemphyxmini0950-forcevectorquantity, def:physics:phyx-mini-0950:phyxminiproblems-problemphyxmini0950-vectorsireadout}
  Cartesian newton readout of the net magnetic force
\end{definition}
\begin{proof}
  This is the declared model definition of force Vector In Newtons.
\end{proof}
\begin{definition}[force Magnitude In Newtons]
  \label{def:physics:phyx-mini-0950:phyxminiproblems-problemphyxmini0950-forcemagnitudeinnewtons}
  \lean{PhyXMiniProblems.ProblemPhyXMini0950.forceMagnitudeInNewtons}
  \uses{def:physics:phyx-mini-0950:phyxminiproblems-problemphyxmini0950-forcevectorquantity, def:physics:phyx-mini-0950:phyxminiproblems-problemphyxmini0950-forcevectorinnewtons}
  Magnitude in newtons of a dimensionful force vector
\end{definition}
\begin{proof}
  This is the declared model definition of force Magnitude In Newtons.
\end{proof}
\begin{definition}[x Axis Index]
  \label{def:physics:phyx-mini-0950:phyxminiproblems-problemphyxmini0950-xaxisindex}
  \lean{PhyXMiniProblems.ProblemPhyXMini0950.xAxisIndex}
  Coordinate 0, the horizontal x-axis shown in the raster
\end{definition}
\begin{proof}
  This is the declared model definition of x Axis Index.
\end{proof}
\begin{definition}[y Axis Index]
  \label{def:physics:phyx-mini-0950:phyxminiproblems-problemphyxmini0950-yaxisindex}
  \lean{PhyXMiniProblems.ProblemPhyXMini0950.yAxisIndex}
  Coordinate 1, the loudspeaker and coil y-axis
\end{definition}
\begin{proof}
  This is the declared model definition of y Axis Index.
\end{proof}
\begin{definition}[z Axis Index]
  \label{def:physics:phyx-mini-0950:phyxminiproblems-problemphyxmini0950-zaxisindex}
  \lean{PhyXMiniProblems.ProblemPhyXMini0950.zAxisIndex}
  Coordinate 2, completing the right-handed Cartesian frame
\end{definition}
\begin{proof}
  This is the declared model definition of z Axis Index.
\end{proof}
\begin{definition}[x Hat]
  \label{def:physics:phyx-mini-0950:phyxminiproblems-problemphyxmini0950-xhat}
  \lean{PhyXMiniProblems.ProblemPhyXMini0950.xHat}
  \uses{def:physics:phyx-mini-0950:phyxminiproblems-problemphyxmini0950-spatialvector, def:physics:phyx-mini-0950:phyxminiproblems-problemphyxmini0950-xaxisindex}
  Positive Cartesian x unit vector
\end{definition}
\begin{proof}
  This is the declared model definition of x Hat.
\end{proof}
\begin{definition}[y Hat]
  \label{def:physics:phyx-mini-0950:phyxminiproblems-problemphyxmini0950-yhat}
  \lean{PhyXMiniProblems.ProblemPhyXMini0950.yHat}
  \uses{def:physics:phyx-mini-0950:phyxminiproblems-problemphyxmini0950-spatialvector, def:physics:phyx-mini-0950:phyxminiproblems-problemphyxmini0950-yaxisindex}
  Positive axial y unit vector
\end{definition}
\begin{proof}
  This is the declared model definition of y Hat.
\end{proof}
\begin{definition}[z Hat]
  \label{def:physics:phyx-mini-0950:phyxminiproblems-problemphyxmini0950-zhat}
  \lean{PhyXMiniProblems.ProblemPhyXMini0950.zHat}
  \uses{def:physics:phyx-mini-0950:phyxminiproblems-problemphyxmini0950-spatialvector, def:physics:phyx-mini-0950:phyxminiproblems-problemphyxmini0950-zaxisindex}
  Positive Cartesian z unit vector
\end{definition}
\begin{proof}
  This is the declared model definition of z Hat.
\end{proof}
\begin{definition}[outward Radial Direction]
  \label{def:physics:phyx-mini-0950:phyxminiproblems-problemphyxmini0950-outwardradialdirection}
  \lean{PhyXMiniProblems.ProblemPhyXMini0950.outwardRadialDirection}
  \uses{def:physics:phyx-mini-0950:phyxminiproblems-problemphyxmini0950-spatialvector, def:physics:phyx-mini-0950:phyxminiproblems-problemphyxmini0950-xhat, def:physics:phyx-mini-0950:phyxminiproblems-problemphyxmini0950-zhat}
  The outward radial unit vector in the coil's x-z plane. The minus sign on the z component fixes the view from positive y: increasing azimuth moves counterclockwise on the viewer's page
\end{definition}
\begin{proof}
  This is the declared model definition of outward Radial Direction.
\end{proof}
\begin{definition}[counterclockwise Tangent Direction]
  \label{def:physics:phyx-mini-0950:phyxminiproblems-problemphyxmini0950-counterclockwisetangentdirection}
  \lean{PhyXMiniProblems.ProblemPhyXMini0950.counterclockwiseTangentDirection}
  \uses{def:physics:phyx-mini-0950:phyxminiproblems-problemphyxmini0950-spatialvector, def:physics:phyx-mini-0950:phyxminiproblems-problemphyxmini0950-xhat, def:physics:phyx-mini-0950:phyxminiproblems-problemphyxmini0950-zhat}
  Unit tangent for counterclockwise current as viewed from positive y
\end{definition}
\begin{proof}
  This is the declared model definition of counterclockwise Tangent Direction.
\end{proof}
\begin{definition}[spatial Cross]
  \label{def:physics:phyx-mini-0950:phyxminiproblems-problemphyxmini0950-spatialcross}
  \lean{PhyXMiniProblems.ProblemPhyXMini0950.spatialCross}
  \uses{def:physics:phyx-mini-0950:phyxminiproblems-problemphyxmini0950-spatialvector}
  Mathlib's cross product transported to Physlib's Euclidean vectors
\end{definition}
\begin{proof}
  This is the declared model definition of spatial Cross.
\end{proof}
\begin{definition}[Diagram Axis]
  \label{def:physics:phyx-mini-0950:phyxminiproblems-problemphyxmini0950-diagramaxis}
  \lean{PhyXMiniProblems.ProblemPhyXMini0950.DiagramAxis}
  The two labelled coordinate axes visible in image 950.png
\end{definition}
\begin{proof}
  This is the declared model definition of Diagram Axis.
\end{proof}
\begin{definition}[Field Arrow Side]
  \label{def:physics:phyx-mini-0950:phyxminiproblems-problemphyxmini0950-fieldarrowside}
  \lean{PhyXMiniProblems.ProblemPhyXMini0950.FieldArrowSide}
  The two symmetric magnetic-field arrows drawn beside the coil
\end{definition}
\begin{proof}
  This is the declared model definition of Field Arrow Side.
\end{proof}
\begin{definition}[Figure Arrow Direction]
  \label{def:physics:phyx-mini-0950:phyxminiproblems-problemphyxmini0950-figurearrowdirection}
  \lean{PhyXMiniProblems.ProblemPhyXMini0950.FigureArrowDirection}
  Qualitative arrow directions used to transcribe the primary raster
\end{definition}
\begin{proof}
  This is the declared model definition of Figure Arrow Direction.
\end{proof}
\begin{definition}[Circulation Sense]
  \label{def:physics:phyx-mini-0950:phyxminiproblems-problemphyxmini0950-circulationsense}
  \lean{PhyXMiniProblems.ProblemPhyXMini0950.CirculationSense}
  Sense of current around the coil when viewed from the positive axis
\end{definition}
\begin{proof}
  This is the declared model definition of Circulation Sense.
\end{proof}
\begin{definition}[expected Field Arrow Direction]
  \label{def:physics:phyx-mini-0950:phyxminiproblems-problemphyxmini0950-expectedfieldarrowdirection}
  \lean{PhyXMiniProblems.ProblemPhyXMini0950.expectedFieldArrowDirection}
  \uses{def:physics:phyx-mini-0950:phyxminiproblems-problemphyxmini0950-fieldarrowside, def:physics:phyx-mini-0950:phyxminiproblems-problemphyxmini0950-figurearrowdirection}
  Expected direction of each blue field arrow in the side view
\end{definition}
\begin{proof}
  This is the declared model definition of expected Field Arrow Direction.
\end{proof}
\begin{definition}[Voice Coil Figure]
  \label{def:physics:phyx-mini-0950:phyxminiproblems-problemphyxmini0950-voicecoilfigure}
  \lean{PhyXMiniProblems.ProblemPhyXMini0950.VoiceCoilFigure}
  \uses{def:physics:phyx-mini-0950:phyxminiproblems-problemphyxmini0950-diagramaxis, def:physics:phyx-mini-0950:phyxminiproblems-problemphyxmini0950-fieldarrowside, def:physics:phyx-mini-0950:phyxminiproblems-problemphyxmini0950-figurearrowdirection}
  Literal and qualitative content of the supplied raster. The real-valued angle labels are figure readouts in degrees, not solved force quantities
\end{definition}
\begin{proof}
  This is the declared model definition of Voice Coil Figure.
\end{proof}
\begin{definition}[Voice Coil Setup]
  \label{def:physics:phyx-mini-0950:phyxminiproblems-problemphyxmini0950-voicecoilsetup}
  \lean{PhyXMiniProblems.ProblemPhyXMini0950.VoiceCoilSetup}
  \uses{def:physics:phyx-mini-0950:phyxminiproblems-problemphyxmini0950-spatialvector, def:physics:phyx-mini-0950:phyxminiproblems-problemphyxmini0950-lengthquantity, def:physics:phyx-mini-0950:phyxminiproblems-problemphyxmini0950-currentmagnitudequantity, def:physics:phyx-mini-0950:phyxminiproblems-problemphyxmini0950-magneticfluxdensitymagnitudequantity, def:physics:phyx-mini-0950:phyxminiproblems-problemphyxmini0950-forcevectorquantity, def:physics:phyx-mini-0950:phyxminiproblems-problemphyxmini0950-diagramaxis, def:physics:phyx-mini-0950:phyxminiproblems-problemphyxmini0950-circulationsense, def:physics:phyx-mini-0950:phyxminiproblems-problemphyxmini0950-voicecoilfigure}
  Independent physical data for the voice coil. In particular, netMagneticForce is not defined from a force formula or answer choice
\end{definition}
\begin{proof}
  This is the declared model definition of Voice Coil Setup.
\end{proof}
\begin{definition}[Matches Voice Coil Scenario]
  \label{def:physics:phyx-mini-0950:phyxminiproblems-problemphyxmini0950-matchesvoicecoilscenario}
  \lean{PhyXMiniProblems.ProblemPhyXMini0950.MatchesVoiceCoilScenario}
  \uses{def:physics:phyx-mini-0950:phyxminiproblems-problemphyxmini0950-voicecoilsetup}
  Prose-level axis and circulation assignments
\end{definition}
\begin{proof}
  This is the declared model definition of Matches Voice Coil Scenario.
\end{proof}
\begin{definition}[Matches Primary Voice Coil Figure]
  \label{def:physics:phyx-mini-0950:phyxminiproblems-problemphyxmini0950-matchesprimaryvoicecoilfigure}
  \lean{PhyXMiniProblems.ProblemPhyXMini0950.MatchesPrimaryVoiceCoilFigure}
  \uses{def:physics:phyx-mini-0950:phyxminiproblems-problemphyxmini0950-expectedfieldarrowdirection, def:physics:phyx-mini-0950:phyxminiproblems-problemphyxmini0950-voicecoilsetup}
  Primary-image evidence: the cylinder and coil, labelled axes, rightward front current arrow, symmetric outward/upward B arrows, and two 60 degree arcs
\end{definition}
\begin{proof}
  This is the declared model definition of Matches Primary Voice Coil Figure.
\end{proof}
\begin{definition}[Matches Problem Readouts]
  \label{def:physics:phyx-mini-0950:phyxminiproblems-problemphyxmini0950-matchesproblemreadouts}
  \lean{PhyXMiniProblems.ProblemPhyXMini0950.MatchesProblemReadouts}
  \uses{def:physics:phyx-mini-0950:phyxminiproblems-problemphyxmini0950-lengthincentimeters, def:physics:phyx-mini-0950:phyxminiproblems-problemphyxmini0950-currentinamperes, def:physics:phyx-mini-0950:phyxminiproblems-problemphyxmini0950-magneticfluxdensityinteslas, def:physics:phyx-mini-0950:phyxminiproblems-problemphyxmini0950-voicecoilsetup}
  Numerical values stated in the prose. The angle is converted from degrees to radians. No force value or answer-choice label occurs in this record
\end{definition}
\begin{proof}
  This is the declared model definition of Matches Problem Readouts.
\end{proof}
\begin{definition}[Has Physical Voice Coil Parameters]
  \label{def:physics:phyx-mini-0950:phyxminiproblems-problemphyxmini0950-hasphysicalvoicecoilparameters}
  \lean{PhyXMiniProblems.ProblemPhyXMini0950.HasPhysicalVoiceCoilParameters}
  \uses{def:physics:phyx-mini-0950:phyxminiproblems-problemphyxmini0950-lengthinmeters, def:physics:phyx-mini-0950:phyxminiproblems-problemphyxmini0950-currentinamperes, def:physics:phyx-mini-0950:phyxminiproblems-problemphyxmini0950-magneticfluxdensityinteslas, def:physics:phyx-mini-0950:phyxminiproblems-problemphyxmini0950-voicecoilsetup}
  Positivity and angular nondegeneracy of the physical setup
\end{definition}
\begin{proof}
  This is the declared model definition of Has Physical Voice Coil Parameters.
\end{proof}
\begin{definition}[Satisfies Circular Voice Coil Geometry]
  \label{def:physics:phyx-mini-0950:phyxminiproblems-problemphyxmini0950-satisfiescircularvoicecoilgeometry}
  \lean{PhyXMiniProblems.ProblemPhyXMini0950.SatisfiesCircularVoiceCoilGeometry}
  \uses{def:physics:phyx-mini-0950:phyxminiproblems-problemphyxmini0950-lengthinmeters, def:physics:phyx-mini-0950:phyxminiproblems-problemphyxmini0950-outwardradialdirection, def:physics:phyx-mini-0950:phyxminiproblems-problemphyxmini0950-counterclockwisetangentdirection, def:physics:phyx-mini-0950:phyxminiproblems-problemphyxmini0950-voicecoilsetup}
  The wire is a circle of the stated radius in the plane normal to y, and its current tangent has the counterclockwise orientation fixed above
\end{definition}
\begin{proof}
  This is the declared model definition of Satisfies Circular Voice Coil Geometry.
\end{proof}
\begin{definition}[Models Axisymmetric Voice Coil Magnetic Field]
  \label{def:physics:phyx-mini-0950:phyxminiproblems-problemphyxmini0950-modelsaxisymmetricvoicecoilmagneticfield}
  \lean{PhyXMiniProblems.ProblemPhyXMini0950.ModelsAxisymmetricVoiceCoilMagneticField}
  \uses{def:physics:phyx-mini-0950:phyxminiproblems-problemphyxmini0950-magneticfluxdensityinteslas, def:physics:phyx-mini-0950:phyxminiproblems-problemphyxmini0950-yhat, def:physics:phyx-mini-0950:phyxminiproblems-problemphyxmini0950-outwardradialdirection, def:physics:phyx-mini-0950:phyxminiproblems-problemphyxmini0950-voicecoilsetup}
  At every point of the wire, the field has the stated constant magnitude. Its axial component points along positive y, while its transverse component is radially outward and makes the stated angle from the axial normal
\end{definition}
\begin{proof}
  This is the declared model definition of Models Axisymmetric Voice Coil Magnetic Field.
\end{proof}
\begin{definition}[Satisfies Voice Coil Magnetic Lorentz Force Law]
  \label{def:physics:phyx-mini-0950:phyxminiproblems-problemphyxmini0950-satisfiesvoicecoilmagneticlorentzforcelaw}
  \lean{PhyXMiniProblems.ProblemPhyXMini0950.SatisfiesVoiceCoilMagneticLorentzForceLaw}
  \uses{def:physics:phyx-mini-0950:phyxminiproblems-problemphyxmini0950-lengthinmeters, def:physics:phyx-mini-0950:phyxminiproblems-problemphyxmini0950-currentinamperes, def:physics:phyx-mini-0950:phyxminiproblems-problemphyxmini0950-forcevectorinnewtons, def:physics:phyx-mini-0950:phyxminiproblems-problemphyxmini0950-spatialcross, def:physics:phyx-mini-0950:phyxminiproblems-problemphyxmini0950-voicecoilsetup}
  The governing magnetic Lorentz law. The interval integral parameterizes one turn by azimuth; diameter / 2 converts its unit tangent to dℓ, and the outer scalar multiplication sums the identical force over all turns. This is a line-integral law in the independently supplied field. It does not state the evaluated axial-force formula requested by the problem
\end{definition}
\begin{proof}
  This is the declared model definition of Satisfies Voice Coil Magnetic Lorentz Force Law.
\end{proof}
\begin{definition}[exact Force Magnitude In Newtons]
  \label{def:physics:phyx-mini-0950:phyxminiproblems-problemphyxmini0950-exactforcemagnitudeinnewtons}
  \lean{PhyXMiniProblems.ProblemPhyXMini0950.exactForceMagnitudeInNewtons}
  \uses{def:physics:phyx-mini-0950:phyxminiproblems-problemphyxmini0950-lengthinmeters, def:physics:phyx-mini-0950:phyxminiproblems-problemphyxmini0950-currentinamperes, def:physics:phyx-mini-0950:phyxminiproblems-problemphyxmini0950-magneticfluxdensityinteslas, def:physics:phyx-mini-0950:phyxminiproblems-problemphyxmini0950-voicecoilsetup}
  The positive magnitude obtained after evaluating the line integral. This helper depends only on input readouts and does not define the independent netMagneticForce field
\end{definition}
\begin{proof}
  This is the declared model definition of exact Force Magnitude In Newtons.
\end{proof}
\begin{definition}[signed Axial Force In Newtons]
  \label{def:physics:phyx-mini-0950:phyxminiproblems-problemphyxmini0950-signedaxialforceinnewtons}
  \lean{PhyXMiniProblems.ProblemPhyXMini0950.signedAxialForceInNewtons}
  \uses{def:physics:phyx-mini-0950:phyxminiproblems-problemphyxmini0950-forcevectorinnewtons, def:physics:phyx-mini-0950:phyxminiproblems-problemphyxmini0950-yaxisindex, def:physics:phyx-mini-0950:phyxminiproblems-problemphyxmini0950-voicecoilsetup}
  The signed component of the net force along the pictured y-axis
\end{definition}
\begin{proof}
  This is the declared model definition of signed Axial Force In Newtons.
\end{proof}
\begin{definition}[Points Along Negative YAxis]
  \label{def:physics:phyx-mini-0950:phyxminiproblems-problemphyxmini0950-pointsalongnegativeyaxis}
  \lean{PhyXMiniProblems.ProblemPhyXMini0950.PointsAlongNegativeYAxis}
  \uses{def:physics:phyx-mini-0950:phyxminiproblems-problemphyxmini0950-spatialvector, def:physics:phyx-mini-0950:phyxminiproblems-problemphyxmini0950-yhat}
  A nonzero force vector points along the negative y-axis
\end{definition}
\begin{proof}
  This is the declared model definition of Points Along Negative YAxis.
\end{proof}
\begin{definition}[Answer Choice]
  \label{def:physics:phyx-mini-0950:phyxminiproblems-problemphyxmini0950-answerchoice}
  \lean{PhyXMiniProblems.ProblemPhyXMini0950.AnswerChoice}
  The four answer labels displayed with the problem
\end{definition}
\begin{proof}
  This is the declared model definition of Answer Choice.
\end{proof}
\begin{definition}[signed Force In Newtons]
  \label{def:physics:phyx-mini-0950:phyxminiproblems-problemphyxmini0950-answerchoice-signedforceinnewtons}
  \lean{PhyXMiniProblems.ProblemPhyXMini0950.AnswerChoice.signedForceInNewtons}
  \uses{def:physics:phyx-mini-0950:phyxminiproblems-problemphyxmini0950-answerchoice}
  Signed axial force in newtons printed beside each answer choice
\end{definition}
\begin{proof}
  This is the declared model definition of signed Force In Newtons.
\end{proof}
\begin{definition}[recorded Answer Choice]
  \label{def:physics:phyx-mini-0950:phyxminiproblems-problemphyxmini0950-recordedanswerchoice}
  \lean{PhyXMiniProblems.ProblemPhyXMini0950.recordedAnswerChoice}
  \uses{def:physics:phyx-mini-0950:phyxminiproblems-problemphyxmini0950-answerchoice}
  Answer label recorded in the supplied dataset metadata
\end{definition}
\begin{proof}
  This is the declared model definition of recorded Answer Choice.
\end{proof}
\begin{definition}[answer Choice Error In Newtons]
  \label{def:physics:phyx-mini-0950:phyxminiproblems-problemphyxmini0950-answerchoiceerrorinnewtons}
  \lean{PhyXMiniProblems.ProblemPhyXMini0950.answerChoiceErrorInNewtons}
  \uses{def:physics:phyx-mini-0950:phyxminiproblems-problemphyxmini0950-voicecoilsetup, def:physics:phyx-mini-0950:phyxminiproblems-problemphyxmini0950-signedaxialforceinnewtons, def:physics:phyx-mini-0950:phyxminiproblems-problemphyxmini0950-answerchoice}
  Absolute error from a displayed signed axial-force value
\end{definition}
\begin{proof}
  This is the declared model definition of answer Choice Error In Newtons.
\end{proof}
\begin{definition}[Matches Displayed Signed Force]
  \label{def:physics:phyx-mini-0950:phyxminiproblems-problemphyxmini0950-matchesdisplayedsignedforce}
  \lean{PhyXMiniProblems.ProblemPhyXMini0950.MatchesDisplayedSignedForce}
  \uses{def:physics:phyx-mini-0950:phyxminiproblems-problemphyxmini0950-voicecoilsetup, def:physics:phyx-mini-0950:phyxminiproblems-problemphyxmini0950-answerchoice, def:physics:phyx-mini-0950:phyxminiproblems-problemphyxmini0950-answerchoiceerrorinnewtons}
  Agreement with a force value printed to the nearest 0.001 N
\end{definition}
\begin{proof}
  This is the declared model definition of Matches Displayed Signed Force.
\end{proof}
\begin{definition}[Is Unique Closest Signed Force Choice]
  \label{def:physics:phyx-mini-0950:phyxminiproblems-problemphyxmini0950-isuniqueclosestsignedforcechoice}
  \lean{PhyXMiniProblems.ProblemPhyXMini0950.IsUniqueClosestSignedForceChoice}
  \uses{def:physics:phyx-mini-0950:phyxminiproblems-problemphyxmini0950-voicecoilsetup, def:physics:phyx-mini-0950:phyxminiproblems-problemphyxmini0950-answerchoice, def:physics:phyx-mini-0950:phyxminiproblems-problemphyxmini0950-answerchoiceerrorinnewtons}
  A displayed signed force is strictly closer than every alternative
\end{definition}
\begin{proof}
  This is the declared model definition of Is Unique Closest Signed Force Choice.
\end{proof}
% --- Archon declaration topology end ---
% --- Archon physics formalization source end ---
